% =====================================================================
%  ch04.tex —— 第 4 章：原书 p.51–57（PDF 60–66）
%  4.1 余切丛 ～ 4.6 接触流形
% =====================================================================

\chapter{各种辛流形}

% [原文 p.51]
辛流形大致可分为完全（exact）的与紧致的两类，二者各有长处与短处。也有既不完全也不紧致的辛流形，但这类对象处理起来更为困难。作为完全辛流形的例子，举出余切丛与 Stein 流形；作为紧致辛流形的例子，举出射影簇。此外，作为与表示论等领域相关的重要辛流形，引入余伴随轨道。然后，介绍作为完全辛流形与紧致辛流形之「恰到好处」折中的 Liouville 区域，最后触及出现在其边界上的接触流形的概念。

\takeaway{完全辛流形有餘切丛与 Stein 流形，紧致者有射影簇；余伴随轨道与 Liouville 区域把二者连起来。}

\section{余切丛}

称辛流形 $(M,\omega)$ 为\textbf{完全的}（exact），是指存在 1 次微分形式 $\theta$ 使 $\omega = \rd\theta$。这个 $\theta$ 称为 \textbf{Liouville 形式}（Liouville 1-form）。关于辛形式与 Liouville 形式互为对偶的向量场 $Z$ 称为 \textbf{Liouville 向量场}（Liouville vector field）：
\begin{equation}
  \iota_Z\omega = \theta.
  \tag{4.1.1}
\end{equation}

完全辛流形中最重要且古典的例子，是微分流形 $N$ 的余切丛 $T^*N$。$T^*N$ 上的 1 次微分形式 $\theta$ 对 $X\in T_{(q,p)}(T^*N)$ 由
\begin{equation}
  \theta(X) = \langle p,\pi_*X\rangle
  \tag{4.1.2}
\end{equation}
定义。这里 $\pi\colon T^*N\to N$ 是自然射影，$q\in N$ 且 $p\in\pi^{-1}(q) = T^*_qN$，(4.1.2) 的右端是 $p\in T^*_qN$ 与 $\pi_*X\in T_pN$ 的配对。$N$ 的局部坐标记为 $(q_1,\dots,q_n)$ 时，此坐标在纤维 $T^*_qN$ 上诱导的坐标记为 $(p_1,\dots,p_n)$，则
% [原文 p.52]
\begin{equation}
  \theta = \sum_{i=1}^{n}p_i\rd q_i
  \tag{4.1.3}
\end{equation}
成立。从而
\begin{equation}
  \rd\theta = \sum_{i=1}^{n}\rd p_i\wedge\rd q_i
  \tag{4.1.4}
\end{equation}
成立，由此表式可知 $\rd\theta$ 是非退化的。又，$\rd\theta$ 是完全形式，故特别是闭形式。

% [Tutor Note: 背景——完全辛流形即 $\omega$ 恰当的辛流形；余切丛上的 Liouville 向量场就是沿纤维方向伸缩的 Euler 向量场 $Z=\sum_i p_i\,\partial_{p_i}$，故称「Liouville」实与 Euler 齐次性同源。]

\takeaway{余切丛上的重言 1 形式给出完全辛流形，其 Liouville 向量场沿纤维方向伸缩。}

\section{Stein 流形}

不限于余切丛的完全辛流形，其重要例子是 Stein 流形。

\begin{definition}{4.2.1}
复流形 $X$ 上的光滑函数 $S\colon X\to\R$，若 Levi 形式（Levi form）$\lambda = \sqrt{-1}\partial\bar\partial S$ 为正定值，则称为\textbf{多重劣调和的}（plurisubharmonic）。
\end{definition}

注意，令 $d^cf := \rd f\circ J$，则 $\sqrt{-1}\partial\bar\partial f = -\rd\rd^cf$。

\begin{definition}{4.2.2}
连续函数 $S\colon X\to\R$，若对任意 $K\in\R$，$S^{-1}\bigl((-\infty,K]\bigr)$ 都是紧致的，则称为\textbf{皆既的}（exhaustive）。
\end{definition}

由于 Levi 问题得到了肯定的解决，「复流形具有皆既的多重劣调和函数」成为其为 \textbf{Stein 流形}（Stein manifold）的必要充分条件。Stein 流形的另一刻画是：它是仿射空间 $\C^m$ 的闭复子流形。

由 Nash--Tognoli 定理（例如参看 [69]），对任意 $n$ 维微分流形 $N$，都存在某个自然数 $l$ 与 $l$ 个 $2n+1$ 元实系数多项式 $f_1,\dots,f_l\in\R[x_1,\dots,x_{2n+1}]$，使得这些多项式定出的 $\R^{2n+1}$ 的代数子集
\begin{equation}
  U_\R := \{\bm{x}\in\R^{2n+1}\mid f_1(\bm{x}) = \dots = f_l(\bm{x}) = 0\}
  \tag{4.2.1}
\end{equation}
与 $N$ 微分同胚。这些定义多项式所定出的 $\C^k$ 的子簇（subvariety）
\begin{equation}
  U := \{\bm{x}\in\C^{2n+1}\mid f_1(\bm{x}) = \dots = f_l(\bm{x}) = 0\}
  \tag{4.2.2}
\end{equation}
一般在 $\C^{2n+1}\setminus\R^{2n+1}$ 上带有奇点，但奇点集合 $U^{\mathrm{sing}}\subset U$ 成为由适当的实多项式 $g_1,\dots,g_m$ 定出的 $\C^{2n+1}$ 的子簇。此时，
\begin{equation}
  M = \{(\bm{x},y)\in U\times\C\mid (g_1(\bm{x})^2+\dots+g_m(\bm{x})^2)y = 1\}
  \tag{4.2.3}
\end{equation}

% [Tutor Note: 校勘——(4.2.2) 之前的正文原书作「$\C^k$ 的子簇」，但 $f_i$ 定义在 $\C^{2n+1}$ 上，同页下文亦作 $\C^{2n+1}$，疑为 $\C^{2n+1}$ 之误；照录原文。]

% [Tutor Note: 背景——Levi 问题问的是「多重劣调和函数的存在性是否刻画 Stein 流形」，由 Oka、Bremermann 等人在 1950 年代给出肯定答案（Oka 定理）。]

% [原文 p.53]
与 $U\setminus U^{\mathrm{sing}}$ 同型，且 $M_\R\cong U_\R\cong N$。这个 $M$ 是 $\C^{2n+1}$ 的闭子流形，故为 Stein 流形，并且以 $N$ 作为 Lagrange 子流形包含在内。因此由定理 3.2.4 可知，$M$ 把 $T^*N$ 作为开子流形包含在内。

\takeaway{Stein 流形即具皆既多重劣调和函数的复流形，可嵌入 $\C^{2n+1}$ 而成为完全辛流形。}

\section{射影簇}

复仿射空间 $\C^{n+1}$ 上自然地作用着代数环面 $\C^\times$：
\begin{equation}
  \C^\times\ni\alpha\colon\bm{x} = (x_1,\dots,x_{n+1})\mapsto\alpha\cdot\bm{x} = (\alpha x_1,\dots,\alpha x_{n+1}).
  \tag{4.3.1}
\end{equation}
此作用在原点 $\bm{0}\in\C^{n+1}$ 之外是自由的，商空间 $(\C^{n+1}\setminus\bm{0})/\C^\times$ 上自然带有复流形的结构。称之为 \textbf{射影空间}（projective space），记作 $\PP^n$。能作为闭子流形全纯地嵌入射影空间的复流形，称为 \textbf{射影簇}（projective manifold）。射影簇是 Stein 流形的紧致类似物。由 Chow 定理可知，射影空间的任意复子流形都是代数簇，即可以表为有限个齐次多项式 $f_1,\dots,f_k\in\C[x_1,\dots,x_{n+1}]$ 的公共零点集。

射影空间至少有 3 种表示，它们各自反映射影空间的 3 个侧面。第 1 种表示是上面举出的仿射空间的环面作用之商
\begin{equation}
  \PP^n = (\C^{n+1}\setminus\{\bm{0}\})/\C^\times,
  \tag{4.3.2}
\end{equation}
它强调射影空间作为环面簇（toric variety）的侧面。关于环面簇，第 8 章将更详细地涉及。

第 2 种表示是复 Lie 群（或代数群）关于抛物子群的商
\begin{equation}
  \PP^n = G_\C/P,
  \tag{4.3.3}
\end{equation}
这里 $G_\C$ 是 $\GL_{n+1}(\C)$，或 $G = \SL_{n+1}(\C)$，而 $P\subset G_\C$ 是由
\begin{equation}
  P = \left\{\begin{pmatrix}A & B\\ 0 & D\end{pmatrix}\in G_\C\;\middle|\;
  A\in\GL_1(\C),\ B\in\operatorname{Mat}_{1\times n}(\C),\ D\in\GL_n(\C)\right\}
  \tag{4.3.4}
\end{equation}
形状的分块上三角矩阵构成的子群。它强调射影空间作为复流形（或代数簇）的侧面。

第 3 种表示是 4.4 节要讲的余伴随轨道
\begin{equation}
  \PP = \mathcal{O}_\lambda
  \tag{4.3.5}
\end{equation}

% [Tutor Note: 校勘——(4.3.3) 之后原书作「$G_\C$ 是 $\GL_{n+1}(\C)$ 或 $G=\SL_{n+1}(\C)$」，第二处记号当为 $G_\C=\SL_{n+1}(\C)$；照录原文。]

% [原文 p.54]
并强调射影空间作为辛流形的侧面。这里 $G = \U(n)$ 是酉群，$\mathfrak{g} = \{\xi\in\operatorname{Mat}_\C(n)\mid\xi^* = -\xi\}$ 是它的 李代数，$\lambda = \operatorname{diag}(\lambda,0,\dots,0)\in\mathfrak{g}$。不过，利用内积
\begin{equation}
  \mathfrak{g}\times\mathfrak{g}\to\R,\qquad (X,Y)\mapsto\tr(XY)
  \tag{4.3.6}
\end{equation}
把 $\mathfrak{g}$ 与 $\mathfrak{g}^*$ 视为同一。此内积在伴随作用下不变：
\begin{align}
  \tr(\ad_Z(X)Y) &= \tr([Z,X]Y)
  \tag{4.3.7}\\
  &= \tr(ZXY)-\tr(XZY)
  \tag{4.3.8}\\
  &= \tr(XYZ)-\tr(XZY)
  \tag{4.3.9}\\
  &= -\tr(X[Z,Y])
  \tag{4.3.10}\\
  &= -\tr(X\ad_Z(Y)),
  \tag{4.3.11}
\end{align}
故由这一同一视，余伴随作用移到伴随作用上（参看 (4.4.7)）。由「歪 Hermite 矩阵乘以 $\sqrt{-1}$ 即成为 Hermite 矩阵」这一点立即可知，歪 Hermite 矩阵可用酉矩阵对角化，故 $\mathcal{O}_\lambda$ 是全体固有值为 $(\lambda,0,\dots,0)$ 的歪 Hermite 矩阵所构成的空间。$(\lambda,0,\dots,0)$ 的固定部分群是对角矩阵，于是得到表示
\begin{equation}
  \mathcal{O}_\lambda\cong G/T.
  \tag{4.3.12}
\end{equation}
正如记号所暗示的，复 Lie 群 $G_\C$ 是实 Lie 群 $G$ 的复化（也就是说，$G_\C$ 的 李代数 $\mathfrak{g}_\C$ 是 $G$ 的 李代数 $\mathfrak{g}$ 的复化）。由于 $G$ 在 $G_\C/P$ 上的作用是推移的，
\begin{equation}
  G_\C/P\cong G/(P\cap G)
  \tag{4.3.13}
\end{equation}
成立。
\begin{equation}
  P\cap G = \left\{\begin{pmatrix}A & 0\\ 0 & D\end{pmatrix}\in G_\C\;\middle|\;
  A\in\U(1),\ D\in\U(n)\right\},
  \tag{4.3.14}
\end{equation}
所以
\begin{equation}
  \PP^n\cong\U(n+1)/\bigl(\U(1)\times\U(n)\bigr)
  \tag{4.3.15}
\end{equation}
成立。

射影空间带有以
\begin{equation}
  \omega_{\mathrm{FS}} = \sqrt{-1}\partial\bar\partial\log\bigl(|z_1|^2+\dots+|z_{n+1}|^2\bigr)
  \tag{4.3.16}
\end{equation}
\begin{equation}
  = \sqrt{-1}\partial\left(\frac{\sum_{i=1}^{n+1}z_i\rd\bar z_i}{|z_1|^2+\dots+|z_{n+1}|^2}\right)
  \tag{4.3.17}
\end{equation}
% [原文 p.55]

\begin{equation}
  = \sqrt{-1}\,\frac{\sum_{i=1}^{n+1}\rd z_i\wedge\rd\bar z_i
    -\Bigl(\sum_{i=1}^{n+1}\bar z_i\wedge\rd z_i\Bigr)\wedge
    \Bigl(\sum_{i=1}^{n+1}z_i\wedge\rd\bar z_i\Bigr)}
    {\bigl(|z_1|^2+\dots+|z_{n+1}|^2\bigr)^2}
  \tag{4.3.18}
\end{equation}
为 Kähler 形式的 Kähler 度量。称之为 \textbf{Fubini--Study 度量}（Fubini--Study metric），$\omega_{\mathrm{FS}}$ 称为 \textbf{Fubini--Study 形式}（Fubini--Study form）。由定理 3.5.2，任意射影簇都是 Kähler 流形，特别是辛流形。

% [Tutor Note: 难点——(4.3.16)$\to$(4.3.18) 是对 $\log\sum|z_i|^2$ 逐次求 $\partial$、$\bar\partial$ 并把 $\bar\partial\partial\log$ 的分母通分的结果；三个式子都带同一个编号序列。]

% [Tutor Note: 校勘——(4.3.18) 原书在 $\bar z_i$ 与 $\rd z_i$ 之间排了 $\wedge$（即 $\bar z_i\wedge\rd z_i$），按 Fubini--Study 形式的通常写法此处应为函数与形式的乘积；照录原文。]

\takeaway{射影空间有三种表示，其 Fubini--Study 形式给出 Kähler 结构，故射影簇是辛流形。}

\section{余伴随轨道}

设 $G$ 为 Lie 群，其 李代数记为 $\mathfrak{g}$。$G$ 在 $G$ 上的伴随作用（adjoint action）
\begin{equation}
  \operatorname{Ad}\colon G\times G\to G,\qquad (g,h)\mapsto ghg^{-1}
  \tag{4.4.1}
\end{equation}
微分后所得的 $G$ 在 $\mathfrak{g}$ 上的作用称为 $G$ 在 $\mathfrak{g}$ 上的\textbf{伴随作用}（adjoint action），并用同一记号
\begin{equation}
  \operatorname{Ad}\colon G\times\mathfrak{g}\to\mathfrak{g}
  \tag{4.4.2}
\end{equation}
表示。此外，把此作用看作 $G$ 在向量空间 $\mathfrak{g}$ 上的线性表示，则称为\textbf{伴随表示}（adjoint representation）。此表示的对偶表示
\begin{equation}
  \operatorname{Ad}^*\colon G\times\mathfrak{g}^*\to\mathfrak{g}^*
  \tag{4.4.3}
\end{equation}
由
\begin{equation}
  \langle\operatorname{Ad}^*_g\xi,X\rangle = \langle\xi,\operatorname{Ad}_{g^{-1}}X\rangle
  \tag{4.4.4}
\end{equation}
定义，称为 $G$ 的\textbf{余伴随表示}（coadjoint representation）。把此表示看作 $G$ 的作用时，称为 $G$ 在 $\mathfrak{g}^*$ 上的\textbf{余伴随作用}（coadjoint action）。取 $\mathfrak{g}$ 的基底把伴随表示用矩阵表示时，$\operatorname{Ad}^*_g$ 关于 $\mathfrak{g}^*$ 的对偶基底的表示矩阵，是 $\operatorname{Ad}_g$ 的表示矩阵的转置之逆矩阵。

把 $G$ 的伴随表示微分，便定义出 $\mathfrak{g}$ 的伴随表示
\begin{equation}
  \ad\colon\mathfrak{g}\times\mathfrak{g}\to\mathfrak{g},\qquad (X,Y)\mapsto\ad_X(Y) = [X,Y].
  \tag{4.4.5}
\end{equation}
把 $G$ 的余伴随表示微分所得的 $\mathfrak{g}$ 的余伴随表示
\begin{equation}
  \ad^*\colon\mathfrak{g}\times\mathfrak{g}^*\to\mathfrak{g}^*
  \tag{4.4.6}
\end{equation}
与 $\mathfrak{g}$ 的伴随表示由
\begin{equation}
  \langle\ad^*_Y(\xi),X\rangle = \langle\xi,-\ad_YX\rangle,\qquad
  X,Y\in\mathfrak{g},\ \xi\in\mathfrak{g}^*
  \tag{4.4.7}
\end{equation}

% [原文 p.56]
相联系。

对 李代数 $\mathfrak{g}$，其双对空间 $\mathfrak{g}^*$ 自然带有 Poisson 结构。对 $f,g\in C^\infty(\mathfrak{g}^*)$ 与 $\xi\in\mathfrak{g}^*$，
\begin{equation}
  T^*_\xi\mathfrak{g}^* = (T_\xi\mathfrak{g}^*)^*\cong(\mathfrak{g}^*)^*\cong\mathfrak{g}
  \tag{4.4.8}
\end{equation}
上带有自然的 Lie 括弧，故用它将 $(\rd f)_\xi,(\rd g)_\xi\in\mathfrak{g}$ 取 Lie 括弧所得的 $\mathfrak{g}$ 的元在 $\xi\in\mathfrak{g}^*$ 处求值，就定义为 $\{f,g\}$ 在 $\xi$ 处的值：
\begin{equation}
  \{f,g\}(\xi) := \langle\xi,[(\rd f)_\xi,(\rd g)_\xi]\rangle.
  \tag{4.4.9}
\end{equation}
$[-,-]$ 的 Jacobi 恒等式导出 $\{-,-\}$ 的 Jacobi 恒等式，$\mathfrak{g}^*$ 由此成为 Poisson 流形。$\{-,-\}$ 称为 $\mathfrak{g}^*$ 的 \textbf{Lie--Poisson 括号}（Lie--Poisson bracket），$\{-,-\}$ 所定出的 $\mathfrak{g}^*$ 的 Poisson 结构称为 $\mathfrak{g}^*$ 的 \textbf{Lie--Poisson 结构}（Lie--Poisson structure）。

\begin{remark}{4.4.1}
若 Lie 群带有 Poisson 结构，且 Lie 群的积成为 Poisson 流形的射，则称为 Poisson--Lie 群，但这与 Lie--Poisson 结构是完全不同的概念。
\end{remark}

下面的主张是到此为止所述诸事的显然归结，但鉴于其重要性，仍作为定理陈述：

\begin{theorem}{4.4.2}
Lie 群 $G$ 的余伴随轨道 $\mathcal{O}\subset\mathfrak{g}^*$ 是 $\mathfrak{g}^*$ 的 Lie--Poisson 结构的辛叶。
\end{theorem}

因此余伴随轨道自然带有辛结构，此时的辛形式称为 \textbf{Kostant--Kirillov 形式}（Kostant--Kirillov form）。

% [Tutor Note: 背景——$\mathfrak{g}^*$ 上的 Lie–Poisson 括号是 Kirillov、Kostant、Souriau 于 1960 年代建立的「轨道方法」的出发点；余伴随轨道上的辛形式即由该括号的辛叶结构诱导。]

\takeaway{$\mathfrak{g}^*$ 的 Lie--Poisson 辛叶即余伴随轨道，其辛形式称 Kostant--Kirillov 形式。}

\section{Liouville 区域}

紧致的带边流形 $X$ 与其上的 1 形式 $\theta$ 的组 $(X,\theta)$，在满足
\begin{itemize}[leftmargin=2em]
  \item $\omega = \rd\theta$ 是辛形式；
  \item 由条件 $\iota_Z\omega = \theta$ 唯一确定的 Liouville 向量场（Liouville vector field）$Z$ 沿 $X$ 的边界指向外侧（处处不与边界相切）
\end{itemize}
这 2 个条件时，称为 \textbf{Liouville 区域}（Liouville domain）。

对任意 $K\in\R$，用皆既的多重劣调和函数 $S\colon X\to\R$ 切出 Stein 流形 $X$ 所得的带边流形 $S^{-1}\bigl((-\infty,K]\bigr)$ 称为 \textbf{Stein 区域}（Stein domain）。组 $\bigl(S^{-1}((-\infty,K]),\ \theta = d^cS\bigr)$ 是 Liouville 区域的重要例子。Liouville 向量场 $Z$ 由 $X$ 的 Kähler 度量下 $S$ 的梯度向量场 $\nabla S$ 给出。

% [Tutor Note: 难点——$Z=\nabla S$ 的验证：$\omega=\rd\rd^cS$，故 $\iota_{\nabla S}\omega = \rd S$ 的两边都在 $\rd\rd^cS$ 的框架下成立；关键是用 Kähler 度量把 1 形式 $\rd S$ 与向量场等同。]

\takeaway{Liouville 区域是带外指 Liouville 向量场的紧致带边辛流形，Stein 区域是其典型例子。}

% [原文 p.57]
\section{接触流形}

设 $n$ 为自然数。$2n+1$ 维流形 $N$ 上的\textbf{接触形式}（contact form）是指微分 1 形式 $\theta$，使
\begin{equation}
  \theta\wedge(\rd\theta)^n
  \tag{4.6.1}
\end{equation}
处处不为 0。$N$ 上的\textbf{接触结构}（contact structure）是指切丛 $TN$ 的部分丛中，能由某个接触形式 $\theta$ 写成 $\xi = \ker\theta$ 者。2 个接触形式 $\theta$ 与 $\theta'$ 定出同一接触结构的必要充分条件是：存在函数 $\phi\colon N\to\R$ 使 $\theta' = e^\phi\theta$。

接触流形是辛流形在奇数维的类似物。构造接触流形的方法之一，是对余切丛 $T^*W$ 与其上显含时间的哈密顿函数 $H\colon T^*W\times\R\to\R$，考虑 $T^*W\times\R$ 上的 Poincaré--Cartan 积分不变量（Poincaré--Cartan integral invariant）
\begin{equation}
  \eta = \bm{p}\cdot\rd\bm{q} - H\rd t.
  \tag{4.6.2}
\end{equation}
另一种构造方法是，对余切丛 $T^*W$ 与其上不显含时间的哈密顿函数 $H\colon T^*W\to\R$，考虑能量超曲面 $H^{-1}(c)\subset T^*W$。无论哪种构造，所得流形成为接触流形的条件都构成对哈密顿函数的非平凡条件。第 1 种构造的典型例是 $H=1$ 的情形，此时的接触形式由 $\eta = \bm{p}\cdot\rd\bm{q}-\rd t$ 给出。第 2 种构造的典型且重要的例子是多重劣调和函数的等位集合，它给出 Liouville 区域的边界。

对接触流形 $(N,\theta)$，$\bigl(N\times\R,\ e^t(\rd\theta+\rd t\wedge\theta)\bigr)$ 是辛流形。称之为 $N$ 的\textbf{辛化}（symplectization）。

保持接触结构的微分同胚称为\textbf{接触同相}（contactomorphism）。微分同胚 $\psi\colon N\to N$ 为接触同相的必要充分条件是：存在函数 $\phi\colon N\to\R$，使 $N$ 的接触形式 $\theta$ 满足 $\psi^*\theta = e^\phi\theta$。

% [Tutor Note: 背景——接触结构是「奇维的辛结构」：$N\times\R$ 上的辛化形式正是把 (4.6.2) 的 $\eta$ 沿 $t$ 方向延拓所得，故接触流形与辛流形的关系一如射影空间与仿射空间的关系。]

\takeaway{接触流形是奇维的辛类比，其辛化 $N\times\R$ 上带有标准辛形式。}

\chapterreview{本章按「完全」与「紧致」两条线索列举辛流形。完全辛流形以余切丛为典范，其 Liouville 形式的重言性使辛形式恰当，Liouville 向量场则沿纤维方向伸缩。Stein 流形由皆既的多重劣调和函数刻画，经 Nash--Tognoli 定理的构造可嵌入 $\C^{2n+1}$ 而成为完全辛流形。紧致方面，射影空间有环面商、抛物子群商与余伴随轨道三种表示，其 Fubini--Study 形式给出 Kähler 结构，故射影簇是辛流形。余伴随轨道作为 $\mathfrak{g}^*$ 上 Lie--Poisson 结构的辛叶，自然带有 Kostant--Kirillov 形式，是表示论与辛几何的交汇点。介于两者之间的 Liouville 区域以 Stein 区域为典型例，其边界上的接触流形则是奇维的辛类比。}
